%%%%%%%%%%%%%%%%%%%%%%%%%%%%%%%%%%%%%%%%%%%%%%%%%%%%%%%%%%%%%%%%%%%%%%%%%%%%%%%
%% Capítulo 1 --- Introdução
%%
%% Citações: abntex2cite (sistema autor-data). \cite{} → (AUTOR, ano);
%%           \citeonline{} → Autor (ano).
%%
%% MARCAÇÃO DE TRECHO EM ANDAMENTO --------------------------------------------
%% \emandamento{...} destaca em vermelho os trechos que mudam com as medições do
%% dispositivo montado. É marcação de rascunho.
%%
%% Para REMOVER todas as marcas na versão final, troque a definição abaixo por:
%%     \newcommand{\emandamento}[1]{#1}
%% ou apague as chamadas \emandamento{...} do texto.
%%%%%%%%%%%%%%%%%%%%%%%%%%%%%%%%%%%%%%%%%%%%%%%%%%%%%%%%%%%%%%%%%%%%%%%%%%%%%%%

\makeatletter
\@ifundefined{emandamento}{%
  \@ifundefined{textcolor}%
    {\newcommand{\emandamento}[1]{#1}}%
    {\newcommand{\emandamento}[1]{\textcolor{red}{#1}}}%
}{}
\makeatother

%% Seção primária numerada: a NBR 14724 (5.2.3) e a NBR 6024 só admitem sem número os
%% títulos que listam, e a Introdução não está entre eles. O modelo do iCEV a deixa sem
%% número; a decisão de seguir a norma é do usuário (2026-11-01).
\chapter{Introdução}
\label{cap:introducao}

A vibração é a grandeza mais usada para acompanhar a condição de máquinas rotativas
\cite{randall2011livro}. A norma ISO 20816-3:2022, da International Organization for
Standardization (ISO), define como avaliá-la em máquinas industriais acima de
\SI{15}{\kilo\watt}, com rotação entre \SI{120}{r/min} e \SI{30000}{r/min}, e pede um
instrumento com resposta plana em \emph{ao menos}
\SIrange{10}{1000}{\hertz}~\cite{iso208163_2022}. Acelerômetros piezoelétricos atendem a
esse requisito, mas eles, seus condicionadores e os analisadores associados custam caro.
Por isso, o monitoramento costuma ficar restrito às máquinas consideradas críticas, e a
maior parte dos motores de uma planta continua sem instrumentação.

Uma alternativa estudada na literatura recente é o nó sensor de baixo custo: uma unidade
de medição inercial (IMU) baseada em sistemas microeletromecânicos (MEMS), ligada a um
microcontrolador que executa a detecção de anomalias e envia o resultado por rádio. O
algoritmo Isolation Forest~\cite{liu2008} já foi treinado e executado dentro de um
ESP32~\cite{antonini2023}, há variantes criadas para reduzir seu custo em aprendizado de
máquina para dispositivos de recursos limitados (TinyML)~\cite{barbariol2022}, e a
combinação de microcontrolador, acelerômetro MEMS de consumo e detecção não supervisionada
já foi aplicada à manutenção preditiva~\cite{kolok2025}. Esses trabalhos executam o detector
no microcontrolador. Eles não indicam quais falhas um sensor desse tipo consegue observar.

Duas etapas determinam o que essa medição entrega. A aquisição define que parte da
vibração chega ao processamento e depende da banda e do piso de ruído do sensor. A
estimação define o que se faz com as amostras: como integrar a aceleração, onde colocar o
limiar do detector, que métrica informar e como separar treino e teste. Uma componente que
não foi adquirida não pode ser recuperada depois.

Sobre a aquisição, os estudos revisados concordam entre si. Nenhum dos estudos de adequação de
sensor revisados testou dispositivo com limite superior abaixo de \SI{1500}{\hertz}~\cite{albarbar2008},
e os grupos que diagnosticam defeito de rolamento com MEMS escolhem sensores de dezenas de
quilohertz~\cite{jakobsen2024,landi2023}, porque a assinatura do defeito aparece modulada
em uma portadora de alta frequência~\cite{antoni2007,randall2011}. Sobre a estimação, há
menos trabalhos. O escore do Isolation Forest usa \num{0,5} como referência de
interpretação, sem calibração para os dados~\cite{liu2008}, e a implementação de
referência transforma esse valor em um deslocamento fixo, que não altera a ordem dos
escores~\cite{scikitlearn2026}. Os trabalhos que relatam bom desempenho com o detector
trocam o valor padrão, mas não testam a alternativa por varredura~\cite{antonini2023,brito2022}.
Entre os trabalhos que executam o detector em microcontrolador, a busca não encontrou
nenhum que variasse o parâmetro de contaminação sobre a mesma base e informasse a fração
da operação normal marcada como anomalia. Também não encontrou taxa de falso alarme
publicada para detector treinado no próprio dispositivo, que é o caso de uma máquina sem
histórico rotulado.

A validação é outro ponto fraco. Janelas consecutivas de um mesmo ensaio têm o mesmo nível
de ruído, o mesmo ganho do sensor e a mesma fixação. Quando essas janelas são sorteadas
entre treino e teste, o modelo pode reconhecer o ensaio em vez da falha.
\citeonline{varejao2025} chamam o fenômeno de \emph{similarity bias} (viés de similaridade)
e relatam que \num{40} de \num{41} estudos revisados sobre uma mesma base pública de
rolamentos usaram delineamento sujeito a ele.

%% PENDENTE (modelo iCEV): motivo da escolha do tema e trajetória do autor. Texto a ser
%% escrito com informações do autor; entra aqui, antes do problema de pesquisa.

\section{Problema de pesquisa e hipóteses}
\label{sec:intro-problema}

O problema de pesquisa deste trabalho é: quais modos de falha de motores uma cadeia de
medição de baixo custo, formada por acelerômetro MEMS de consumo amostrado a
\SI{1}{\kilo\hertz}, microcontrolador ESP32 e enlace de rádio de longo alcance (LoRa),
permite detectar, e quanto as decisões de estimação alteram o desempenho reportado sobre
essa cadeia? O objeto de estudo é o Kaelix, um dispositivo com esses componentes, que
calcula indicadores de vibração no próprio microcontrolador e envia o valor eficaz e o estado
da máquina por LoRa a um receptor.

A hipótese geral é que uma cadeia desse tipo permite detectar desbalanceamento sob uma taxa
de falso alarme fixada, mas não detectar defeito de rolamento sem desbalanceamento
sobreposto, e que as decisões de estimação alteram o desempenho reportado. Ela se desdobra em seis hipóteses, cada uma com a observação que a refutaria.

\textbf{H1, banda.} Com amostragem a \SI{1}{\kilo\hertz} e o filtro interno do sensor, a
        cadeia preserva as harmônicas da rotação, que caracterizam o desbalanceamento, e
        perde a banda de ressonância em que aparece a assinatura de defeito de rolamento.
        Sob a mesma taxa de falso alarme, a detecção de desbalanceamento deve superar a de
        defeito de rolamento sem massa de desbalanceamento sobreposta. A hipótese é
        refutada se a detecção de rolamento for igual ou maior.

\textbf{H2, integração.} A integração no domínio da frequência, restrita à banda exigida
        pela norma, recupera a velocidade eficaz com erro menor que a integração no
        domínio do tempo, com e sem viés contido na especificação do sensor. A hipótese é
        refutada se a integração no tempo tiver erro igual ou menor em alguma dessas
        condições.

\textbf{H3, limiar.} Com treino apenas em operação normal, o limiar padrão do Isolation
        Forest marca como anomalia uma fração da operação normal maior que a marcada por
        um limiar calibrado, sem alterar a ordem dos escores. A hipótese é refutada se
        essa fração for igual ou menor, ou se a ordem mudar.

\textbf{H4, métrica.} Entre a cadeia de referência e a decimada a \SI{1}{\kilo\hertz}, a área parcial sob
        a curva característica de operação do receptor (ROC), restrita à região de falso
        alarme baixo, cai proporcionalmente mais que a área global. A hipótese é refutada
        se a queda relativa da área parcial for igual ou menor.

\textbf{H5, protocolo.} O mesmo modelo, sobre os mesmos dados e descritores, atinge
        acurácia maior sob partição aleatória por janela que sob partição por ensaio. A
        hipótese é refutada se a acurácia for igual ou menor.

\textbf{H6, dispositivo.} \emandamento{O dispositivo montado adquire janelas a
        \SI{1}{\kilo\hertz} e entrega as leituras ao receptor por LoRa com integridade
        verificada. A hipótese é refutada se as janelas forem descartadas pelo limite de
        \SI{2}{\percent} de desvio de taxa do \emph{firmware} ou se o receptor não validar
        os pacotes pela verificação cíclica de redundância (CRC).}

As taxas de detecção por modo de falha e a comparação entre limiares vêm de dado real: gravações
de uma bancada instrumentada, adquiridas a \SI{50}{\kilo\hertz} e decimadas em \emph{software}
para a taxa do dispositivo. Os resultados sobre banda, integração, métrica e protocolo vêm de sinal
sintetizado e não indicam o desempenho em uma máquina real. \emandamento{Os registros do dispositivo em
bancada cobrem aquisição e transmissão e não incluem motor nem falha induzida. As novas
medições com o dispositivo montado ainda estão em andamento.}

\section{Objetivos}
\label{sec:intro-objetivos}

O objetivo geral é determinar quais modos de falha a cadeia de medição do Kaelix permite
detectar e calibrar o limiar do detector para uma taxa de falso alarme fixada.

Os objetivos específicos são:

\begin{alineas}
    \item quantificar a perda de contraste do indicador de defeito de rolamento entre a
          cadeia de referência e a decimada a \SI{1}{\kilo\hertz}, e verificar se o delineamento
          permite atribuir essa perda à banda ou ao método de demodulação;

    \item medir a detecção por modo de falha em dado real decimado para a taxa do
          dispositivo, sob taxa de falso alarme fixada;

    \item determinar o método de integração de aceleração para velocidade compatível com a
          norma, e quantificar o erro dos métodos alternativos contra um sinal cuja
          integral é conhecida analiticamente~\cite{brandt2014,iso208163_2022};

    \item medir a taxa de falso alarme produzida pelo limiar padrão do Isolation Forest e
          por um limiar calibrado, e o efeito do parâmetro de contaminação sobre ela;

    \item comparar a queda de desempenho medida pela área sob a curva ROC global e pela
          área parcial restrita à região de falso alarme baixo;

    \item comparar o desempenho de um classificador sob partição aleatória por janela e sob
          partição por ensaio;

    \item \emandamento{verificar, no dispositivo montado, a aquisição a
          \SI{1}{\kilo\hertz} e a entrega das leituras por LoRa ao receptor.}
\end{alineas}

\section{Justificativa}
\label{sec:intro-justificativa}

A maior parte dos motores fica sem monitoramento por causa do custo. Na manutenção corretiva, a falha ocorre em momento não escolhido e pode danificar componentes
íntegros. Na preventiva, componentes em bom estado são trocados sem necessidade, e falhas
anteriores ao intervalo continuam sem detecção \cite{mobley2002}. A manutenção preditiva
decide a intervenção pelo estado medido do equipamento, mas o custo de instrumentação a
limita às máquinas críticas. Um nó de baixo custo pode atender os demais motores, e quem
for instalá-lo precisa saber quais falhas ele detecta e com que taxa de falso alarme.

Os resultados publicados são difíceis de comparar. Os trabalhos que propõem esses nós
raramente informam a banda e o piso de ruído do sensor, e \num{40} de \num{41} estudos
revisados por \citeonline{varejao2025} usaram delineamento sujeito a vazamento entre treino
e teste. Nessas condições, o leitor não sabe se o número mede a falha, a montagem ou o
protocolo. Este trabalho mede o efeito dessas decisões. Quem precisa detectar defeito de
rolamento pode usar os resultados para escolher o sensor.

Boa parte das decisões que alteram o desempenho medido é tomada no \emph{software}, o que
liga o tema à Engenharia de Software: o método de integração, o limiar, o
protocolo de validação e a correspondência entre o modelo treinado e o executado no
dispositivo. No Kaelix, as funções de processamento são separadas do acesso ao
\emph{hardware} e testadas no computador, sem o dispositivo. Um teste compara o escore do
detector em C++ com o da biblioteca usada no treino, e os erros são propagados até o
pacote enviado ao receptor (Apêndice~\ref{ape:praticas}).

\section{Classificação da pesquisa}
\label{sec:intro-classificacao}

Segundo os critérios de \citeonline{prodanov2013}, esta pesquisa é aplicada quanto à
natureza, explicativa quanto aos objetivos, quantitativa quanto à abordagem e
bibliográfica e experimental quanto aos procedimentos técnicos. Ela é aplicada porque gera
conhecimento dirigido à solução de um problema específico \cite[p.~51]{prodanov2013}, o de
saber o que uma cadeia de baixo custo permite medir em motores. É explicativa porque
procura identificar os fatores que determinam o fenômeno estudado
\cite[p.~53]{prodanov2013}, neste caso o desempenho reportado, em função da banda, do
método de integração, do limiar, da métrica e do protocolo de validação. A abordagem é
quantitativa porque os resultados são expressos em números e analisados com recursos
estatísticos como percentagem, média e desvio-padrão \cite[p.~69]{prodanov2013}.

Os procedimentos técnicos são dois. A revisão da literatura é pesquisa bibliográfica
\cite[p.~54]{prodanov2013}. As comparações de integração, limiar, métrica e protocolo são
pesquisa experimental, porque manipulam uma variável e mantêm as demais sob controle para
observar seu efeito \cite[p.~57]{prodanov2013}: a mesma senoide passa por cinco variantes
de integração, e o mesmo modelo, sobre os mesmos dados, passa por duas formas de partição.
Esses experimentos são computacionais, sobre sinal sintetizado e dado público de bancada.
\emandamento{As medições com o dispositivo montado, ainda em andamento, acrescentam o
experimento físico.}

\section{Etapas do trabalho}
\label{sec:intro-etapas}

A integração foi testada contra um resultado analítico, de modo que um erro de implementação
aparece como desvio. Os demais experimentos compararam cadeias ou métodos sobre os mesmos dados,
com a classe de cada janela conhecida.

O trabalho teve quatro etapas. A primeira foi a revisão da literatura, com o registro das
lacunas encontradas. A segunda foi o estudo da aquisição: um sinal sintetizado passou pela
cadeia de referência e pelas cadeias decimadas a \SI{1}{\kilo\hertz}, e os descritores foram
comparados. Em paralelo, gravações de bancada adquiridas em alta taxa foram decimadas para
medir a detecção por modo de falha à taxa do dispositivo. A terceira foi o estudo da
estimação, com a varredura do parâmetro de contaminação, a avaliação no ponto de operação
de baixo falso alarme e a comparação entre as duas formas de partição. A quarta foi o
desenvolvimento do \emph{firmware} e do receptor e o ensaio do dispositivo em bancada.

Os dados intermediários e o código que gera os resultados estão no repositório do projeto.

\section{Organização do trabalho}
\label{sec:intro-organizacao}

O restante do trabalho está organizado em quatro seções. A
seção~\ref{cap:fundamentacao} apresenta os conceitos usados no texto e os trabalhos
relacionados: vibração de máquinas rotativas e assinaturas de falha, aquisição de sinais,
avaliação pela norma, detecção de anomalias e sua avaliação, nós sensores e enlace LoRa, e
termina com as lacunas encontradas na literatura. A
seção~\ref{cap:desenvolvimento} descreve os experimentos, o dispositivo, o \emph{firmware}
e o receptor. A seção~\ref{cap:resultados} apresenta os resultados e os compara com a
literatura. A seção~\ref{cap:conclusoes} avalia as hipóteses, indica os limites dos
resultados e lista o trabalho futuro, que inclui o projeto de \emph{hardware} próprio.
